\setcounter{section}{0}
\begin{center}
	\zihao{3}
	\textbf{河北工程大学2024$\sim*$2025学年第二学期期末考试(A)卷}
\end{center}

\section{单项选择题}

\begin{enumerate}
	\item 方程 $y''+2x^3y^5+y^6=x^2$ 是\underline{\hspace{2cm}}阶微分方程。
	\begin{choices}
		\item 2
		\item 3
		\item 5
		\item 6
	\end{choices}
	
	\item 已知区域 $D=\{(x,y)\mid x^2+y^2\le 1\}$，则
	$\dfrac{1}{\pi}\iint_D\sqrt{1-x^2-y^2}\,dxdy=$\underline{\hspace{2cm}}。
	\begin{choices}
		\item $\dfrac{2}{3}$
		\item $1$
		\item $\dfrac{4}{3}$
		\item $\dfrac{5}{3}$
	\end{choices}
	
	\item 若级数 $\sum_{n=1}^{\infty}(a_n+b_n)$ 收敛，则级数 $\sum_{n=1}^{\infty}a_n$ 与 $\sum_{n=1}^{\infty}b_n$\underline{\hspace{2cm}}。
	\begin{choices}
		\item 同时收敛
		\item 同时发散
		\item 敛散性不同
		\item 敛散性相同
	\end{choices}
	
	\item 设 $L$ 为直线段 $x=x_0,\ 0\le y\le 2$，则曲线积分 $\int_L 4\,ds=$\underline{\hspace{2cm}}。
	\begin{choices}
		\item $4x_0$
		\item $4$
		\item $8x_0$
		\item $8$
	\end{choices}
	
	\item 若将 $f(x)=x^2\ (0\le x\le \pi)$ 展开为正弦级数，则此级数在 $x=\pi$ 处收敛于\underline{\hspace{2cm}}。
	\begin{choices}
		\item $1$
		\item $0$
		\item $\pi^2$
		\item $-\pi^2$
	\end{choices}
\end{enumerate}

\section{填空题}

\begin{enumerate}
	\item 极限 $\lim\limits_{(x,y)\to(0,0)}\dfrac{5xy}{\sqrt{xy+1}+1}=$\underline{\hspace{2cm}}。
	
	\item 将 $\int_{0}^{1}dx\int_{x}^{\sqrt{3}x}f(x,y)\,dy$ 转化为极坐标下的二次积分为\underline{\hspace{2cm}}。
	
	\item 设 $z=f(x^2-y^2,\ xy)$，且 $f$ 具有一阶连续偏导数，则
	$\dfrac{\partial z}{\partial x}=$\underline{\hspace{2cm}}。
	
	\item 设曲线 $L$ 为圆周 $(x-1)^2+y^2=2$，$L$ 的方向为逆时针方向，则
	$\oint_L y\,dx-x\,dy=$\underline{\hspace{2cm}}。
	
	\item 设区域 $\Omega=\{(x,y,z)\mid x^2+y^2+z^2\le 4\}$，则
	$\iiint_{\Omega}y e^{-(x^2+y^2+z^2)}\,dxdydz=$\underline{\hspace{2cm}}。
\end{enumerate}

\section{计算题}

\begin{enumerate}
	\item 求函数 $u=xyz$ 在 $(1,1,2)$ 处的全微分。
	
	\item 求曲面 $z=x^2+y^2-4$ 在点 $(2,1,1)$ 处的法线方程。
	
	\item 求函数 $f(x,y)=4(x-y)-x^2-y^2$ 的极值。
	
	\item 判定级数 $\sum_{n=1}^{\infty}(-1)^n\dfrac{2}{3^n}$ 的敛散性；如果收敛，是绝对收敛还是条件收敛？
	
	\item 计算二重积分 $\iint_D(3x+2y)\,d\sigma$，其中 $D$ 是由 $x=2,\ y=x,\ y=2x$ 所围成。
	
	\item 设 $L$ 是抛物线 $y=x^2$ 上由点 $A(2,4)$ 到点 $B(-2,4)$ 的一段弧，计算
	$\int_L 2xy\,dx+x^2\,dy$。
\end{enumerate}

\section{解答题}


\begin{enumerate}

\item 已知微分方程 $y'+y=f(x)$，求：

\begin{enumerate}
	\item[(1)] 当 $f(x)=0$ 时，微分方程的通解；
	\item[(2)] 当 $f(x)=e^{-x}$ 时，微分方程的通解。
\end{enumerate}

\item 求幂级数 $\sum_{n=1}^{\infty}n x^{n-1}$ 的收敛域及和函数。

\item 计算曲面积分 $\iint_{\Sigma} z\,dxdy$，其中 $\Sigma$ 是由抛物面 $z=x^2+y^2$ 及平面 $z=2$ 所围成立体的外表面外侧。

\item 设函数 $f(u)$ 具有二阶连续导数，$z=f(e^x\cos y)$ 满足

\begin{align*}
	\frac{\partial^2 z}{\partial x^2}+\frac{\partial^2 z}{\partial y^2}
	=(4z+e^x\cos y)e^{2x},
\end{align*}

若 $f(0)=0,\ f'(0)=0$，求 $f(u)$ 的表达式。


\end{enumerate}
\newpage
